% Acronyms
\newacronym{api}{API}{Application Program Interface}
\newacronym{sdk}{SDK}{Software Development Kit}
\newacronym{uml}{UML}{Unified Modeling Language}
\newacronym{tsa}{TSA}{Termine solo acronimo}

% Glossary
\newglossaryentry{apig}{
    name={API},
    text={Application Program Interface},
    sort=api,
    description={In informatica, un'API è un insieme di procedure a disposizione dei programmatori, di solito raggruppate in un set di strumenti per un compito specifico all'interno di un programma. Il suo scopo è fornire un'astrazione, di solito tra hardware e programmatore o tra software di basso e di alto livello, semplificando l'attività di programmazione}
}

\newglossaryentry{sdkg}{
    name={SDK},
    text={Software Development Kit},
    sort=sdk,
    description={Un Software Development Kit (SDK) è un insieme di strumenti di sviluppo racchiusi in un unico pacchetto installabile, che facilita la creazione di applicazioni mettendo a disposizione un compilatore, un debugger e, talvolta, un framework software. Gli SDK sono in genere specifici per una combinazione di piattaforma hardware e sistema operativo. Molti sviluppatori li usano per abilitare funzionalità avanzate come pubblicità, notifiche push e simili}
}

\newglossaryentry{umlg}{
    name={UML},
    text={Unified Modeling Language},
    sort=uml,
    description={In ingegneria del software, l'Unified Modeling Language (UML) è un linguaggio di modellazione e di specifica basato sul paradigma orientato agli oggetti. L'UML è la «lingua franca» della progettazione e della programmazione a oggetti: gran parte della letteratura di settore lo usa per descrivere soluzioni di analisi e di progettazione in modo sintetico e comprensibile a un pubblico ampio}
}

\newglossaryentry{TermineSenzaAcronimo}{
    name={Nome del termine},
    sort=termine senza acronimo,
    description={Descrizione}
}